%==============================================================
%  lecons/arithmetique/cyclotomiques.tex
%  Les polynômes cyclotomiques
%==============================================================

\lecontitle{Les polynômes cyclotomiques}{Arithmétique — Racines de l'unité et irréductibilité}

%=== NOTATION LOCALE =========================================
\providecommand{\Un}{\mathbb{U}}

%=============================================================
%  § 1 — DÉFINITION ET RÉSULTATS FONDAMENTAUX
%=============================================================
\secnum{navyblue}{1}{Définition et résultats fondamentaux}

\begin{tcolorbox}[thmbox]
  \textbf{Définition.}\enspace\index{polynômes cyclotomiques}%
  \index{racines de l'unité}
  Une racine $\zeta\in\mathbb{C}$ de $X^n-1$ est appelée \emph{racine
  $n$-ième de l'unité}~; elle est dite \emph{primitive} si son ordre dans le
  groupe $\Un_n=\{z\in\mathbb{C}\mid z^n=1\}$ est exactement~$n$, c'est-à-dire
  si $\zeta=e^{2i\pi k/n}$ avec $\pgcd(k,n)=1$. Le \emph{$n$-ième polynôme
  cyclotomique} est le polynôme unitaire dont les racines sont précisément les
  racines primitives $n$-ièmes de l'unité~:
  \[
    \Phi_n(X) \;=\; \prod_{\substack{1\leq k\leq n\\ \pgcd(k,n)=1}}
      \bigl(X-e^{2i\pi k/n}\bigr).
  \]

  \medskip
  \textbf{Théorème (propriétés fondamentales).}
  \begin{itemize}
    \item \textit{Degré.}\enspace\index{indicatrice d'Euler}
      $\deg\Phi_n=\varphi(n)$, où $\varphi$ désigne l'\emph{indicatrice
      d'Euler} (nombre d'entiers de $\{1,\ldots,n\}$ premiers à~$n$).
    \item \textit{Relation fondamentale.}\enspace La partition des racines
      $n$-ièmes de l'unité selon leur ordre $d\mid n$ donne
      \[
        X^n-1 \;=\; \prod_{d\mid n}\Phi_d(X).
      \]
    \item \textit{Intégralité.}\enspace $\Phi_n\in\mathbb{Z}[X]$ et $\Phi_n$
      est unitaire (conséquence de la relation précédente, par récurrence et
      division euclidienne dans $\mathbb{Z}[X]$).
    \item \textit{Irréductibilité.}\enspace
      $\Phi_n$ est \textbf{irréductible} sur $\mathbb{Q}$, et donc (par le
      lemme de Gauss) sur $\mathbb{Z}$. Ainsi $\Phi_n$ est le \emph{polynôme
      minimal} sur $\mathbb{Q}$ de toute racine primitive $n$-ième de l'unité.
    \item \textit{Formule de Möbius.}\enspace\index{fonction de Möbius}
      Par inversion de la relation fondamentale,
      \[
        \Phi_n(X) \;=\; \prod_{d\mid n}\bigl(X^d-1\bigr)^{\mu(n/d)},
      \]
      où $\mu$ est la \emph{fonction de Möbius}.
  \end{itemize}

  \medskip
  \textbf{Cas particuliers.}
  \[
    \Phi_1=X-1,\qquad \Phi_2=X+1,\qquad
    \Phi_p(X)=1+X+\cdots+X^{p-1}=\frac{X^p-1}{X-1}\ \ (p\text{ premier}).
  \]
\end{tcolorbox}

\begin{significationintuitive}
Le polynôme $X^n-1$ se factorise sur $\mathbb{C}$ en facteurs linéaires
distincts $(X-\zeta)$, un par racine $n$-ième de l'unité. Les $\Phi_d$,
$d\mid n$, \emph{regroupent} ces racines selon leur \emph{ordre exact}~:
$\Phi_d$ rassemble exactement les racines d'ordre~$d$. Les polynômes
cyclotomiques sont ainsi les briques indécomposables (sur $\mathbb{Q}$) de
$X^n-1$, et $\Phi_n$ est le « bon » objet algébrique attaché à une racine
primitive $\zeta_n=e^{2i\pi/n}$~: son
irréductibilité fait de $\mathbb{Q}(\zeta_n)$ une extension de degré
$\varphi(n)$, fondement de la théorie des \emph{extensions cyclotomiques}.
\end{significationintuitive}

\decsep{}

%=============================================================
%  § 2 — DÉMONSTRATION
%=============================================================
\secnum{navyblue}{2}{Démonstration}

\begin{ideepreuve}
La relation $X^n-1=\prod_{d\mid n}\Phi_d$ traduit simplement le fait que
chaque racine $n$-ième de l'unité possède un \emph{unique} ordre $d$ divisant
$n$. On en déduit, par récurrence sur $n$ et division euclidienne par un
polynôme \emph{unitaire à coefficients entiers}, que les $\Phi_n$ sont eux-mêmes
à coefficients entiers. La formule de Möbius s'obtient par inversion
multiplicative de cette relation. L'irréductibilité, plus délicate, repose sur
un argument de réduction modulo un premier $p$~: si $\zeta$ est racine du
facteur minimal $f$ de $\Phi_n$, on montre que $\zeta^p$ l'est aussi pour tout
premier $p\nmid n$, d'où $f$ contient toutes les racines primitives.
\end{ideepreuve}

\vspace{10pt}
\rubriq{Preuve}

\begin{mdframed}[style=proofbar]
  \textit{Étape 1 — Relation fondamentale.}
  Toute racine $n$-ième de l'unité $\zeta$ engendre un sous-groupe cyclique de
  $\Un_n$~; son ordre $d=\mathrm{ord}(\zeta)$ divise $n$ (théorème de
  Lagrange), et $\zeta$ est alors racine \emph{primitive} $d$-ième de l'unité.
  Réciproquement, toute racine primitive $d$-ième pour $d\mid n$ vérifie
  $\zeta^n=1$. Comme $X^n-1$ a $n$ racines simples (sa dérivée $nX^{n-1}$ ne
  s'annule qu'en~$0$, qui n'est pas racine de $X^n-1$), on partitionne
  l'ensemble de ses racines selon leur ordre $d\mid n$~:
  \[
    X^n-1 \;=\; \prod_{\zeta^n=1}(X-\zeta)
            \;=\; \prod_{d\mid n}\ \prod_{\mathrm{ord}(\zeta)=d}(X-\zeta)
            \;=\; \prod_{d\mid n}\Phi_d(X).
  \]
  En comparant les degrés, on retrouve au passage l'identité d'Euler
  $\sum_{d\mid n}\varphi(d)=n$, et $\deg\Phi_n=\varphi(n)$.

  \medskip
  \textit{Étape 2 — Intégralité (récurrence \& division euclidienne).}
  On montre $\Phi_n\in\mathbb{Z}[X]$ unitaire par récurrence forte sur~$n$.
  Pour $n=1$, $\Phi_1=X-1$ convient. Pour $n\geq 2$, posons
  $g=\prod_{d\mid n,\ d<n}\Phi_d$. Par hypothèse de récurrence, $g$ est unitaire
  à coefficients entiers, et la relation de l'étape~1 s'écrit $X^n-1=g\,\Phi_n$.
  On invoque le lemme suivant.

  \smallskip\noindent
  \textbf{Lemme (division par un unitaire entier).}\enspace
  \textit{Si $A\in\mathbb{Z}[X]$ et $B\in\mathbb{Z}[X]$ est unitaire, le quotient
  et le reste de la division euclidienne de $A$ par $B$ (dans $\mathbb{Q}[X]$)
  sont à coefficients entiers.}

  \smallskip\noindent
  En effet, l'algorithme de division n'effectue que des soustractions et des
  multiplications par les coefficients de $A$ et $B$~; le coefficient dominant
  de $B$ valant~$1$, aucune division non entière n'intervient. Appliqué à
  $A=X^n-1$ et $B=g$, ce lemme donne $\Phi_n=(X^n-1)/g\in\mathbb{Z}[X]$, et
  $\Phi_n$ est unitaire car $X^n-1$ et $g$ le sont.

  \medskip
  \textit{Étape 3 — Formule de Möbius.}
  La relation $X^n-1=\prod_{d\mid n}\Phi_d$ est une identité multiplicative
  dans le groupe abélien $\mathbb{Q}(X)^\times$ des fractions rationnelles non
  nulles. Or la \emph{formule d'inversion de Möbius} vaut dans tout groupe
  abélien noté multiplicativement~: si $F(n)=\prod_{d\mid n}G(d)$ pour tout
  $n\geq 1$, alors $G(n)=\prod_{d\mid n}F(d)^{\mu(n/d)}$. Appliquée à
  $F(n)=X^n-1$ et $G(d)=\Phi_d$, elle donne
  \[
    \Phi_n(X)=\prod_{d\mid n}\bigl(X^d-1\bigr)^{\mu(n/d)}.
  \]
  (La fonction $\mu$ vérifie $\sum_{d\mid n}\mu(d)=\mathbf{1}_{n=1}$, ce qui
  fonde l'inversion.)

  \medskip
  \textit{Étape 4 — Irréductibilité sur $\mathbb{Q}$ (esquisse).}
  Soit $\zeta=e^{2i\pi/n}$ une racine primitive et $f\in\mathbb{Z}[X]$ son
  polynôme minimal unitaire~; $f$ divise $\Phi_n$ dans $\mathbb{Z}[X]$ (lemme de
  Gauss), écrivons $\Phi_n=f\,g$ avec $g\in\mathbb{Z}[X]$ unitaire. Fixons un
  premier $p\nmid n$. On montre que \emph{$\zeta^p$ est encore racine de~$f$}.

  Sinon, $\zeta^p$ serait racine de~$g$ (car c'est une racine primitive
  $n$-ième, donc de $\Phi_n=fg$), donc $\zeta$ serait racine de $g(X^p)$~; comme
  $f$ est le polynôme minimal de $\zeta$, on aurait $f\mid g(X^p)$ dans
  $\mathbb{Z}[X]$. Réduisons modulo~$p$ dans $\mathbb{F}_p[X]$~: par le
  morphisme de Frobenius, $\overline{g(X^p)}=\overline{g}(X)^p$, donc
  $\overline f\mid \overline g^{\,p}$. Tout facteur irréductible de $\overline f$
  divise alors $\overline g$, si bien que $\overline f$ et $\overline g$ ont un
  facteur commun, et $\overline{\Phi_n}=\overline f\,\overline g$ a une racine
  \emph{double} dans une clôture de $\mathbb{F}_p$. Or $\overline{\Phi_n}$
  divise $X^n-1$, dont la dérivée $nX^{n-1}$ est non nulle modulo~$p$ (car
  $p\nmid n$) et première avec $X^n-1$~: $X^n-1$ est \emph{séparable} sur
  $\mathbb{F}_p$, donc sans racine multiple. Contradiction. Ainsi $\zeta^p$ est
  racine de~$f$.

  Soit enfin $k$ premier à~$n$, que l'on écrit $k=p_1\cdots p_r$ (facteurs
  premiers comptés avec multiplicité)~; aucun $p_i$ ne divise~$n$. L'étape
  précédente vaut en fait pour \emph{toute} racine primitive $n$-ième racine
  de~$f$ (car $f$, irréductible en tant que polynôme minimal, est aussi son
  polynôme minimal)~: appliquée successivement à $\zeta$, $\zeta^{p_1}$,
  $\zeta^{p_1p_2}$, \ldots, elle montre que $\zeta^k$ est racine de~$f$ pour
  tout $k$ avec $\pgcd(k,n)=1$. Mais ces $\zeta^k$ décrivent
  \emph{toutes} les racines primitives $n$-ièmes~; $f$ a donc $\varphi(n)$
  racines, d'où $f=\Phi_n$ et $g=1$. Le polynôme $\Phi_n$ est irréductible sur
  $\mathbb{Q}$.
  \hfill$\blacksquare$
\end{mdframed}

\decsep{}

%=============================================================
%  § 3 — EXEMPLES, CONTRE-EXEMPLES, EXERCICES
%=============================================================
\secnum{exgreen}{3}{Exemples, contre-exemples et exercices}

\rubriq{Exemples}

\begin{mdframed}[style=exbar]
  \textbf{1. Les premiers polynômes cyclotomiques.}
  \[
    \begin{array}{llll}
      \Phi_1=X-1, & \Phi_2=X+1, & \Phi_3=X^2+X+1, & \Phi_4=X^2+1,\\[2pt]
      \Phi_5=\dfrac{X^5-1}{X-1}, & \Phi_6=X^2-X+1, & \Phi_7=\dfrac{X^7-1}{X-1},
        & \Phi_8=X^4+1,\\[6pt]
      \Phi_9=X^6+X^3+1, & \Phi_{10}=X^4-X^3+X^2-X+1, & \Phi_{12}=X^4-X^2+1. &
    \end{array}
  \]
  On vérifie $\deg\Phi_n=\varphi(n)$~: par exemple $\varphi(12)=4=\deg\Phi_{12}$,
  $\varphi(8)=4$, $\varphi(6)=2$. Le calcul de $\Phi_{12}$ illustre Möbius~:
  $\Phi_{12}=\dfrac{(X^{12}-1)(X^2-1)}{(X^6-1)(X^4-1)}=X^4-X^2+1$.

  \smallskip\noindent
  Pour une puissance d'un nombre premier, on dispose de la règle
  $\Phi_{p^k}(X)=\Phi_p\bigl(X^{p^{k-1}}\bigr)$ (une racine $p^k$-ième de
  l'unité est primitive si et seulement si elle n'est pas racine
  $p^{k-1}$-ième)~: ainsi $\Phi_8=\Phi_2(X^4)=X^4+1$ et
  $\Phi_9=\Phi_3(X^3)=X^6+X^3+1$.

  \smallskip\noindent
  \textbf{2. Irréductibilité de $\Phi_p$ par Eisenstein.}\enspace%
  \index{critère d'Eisenstein}
  Pour $p$ premier, $\Phi_p(X)=1+X+\cdots+X^{p-1}$. La translation $X\mapsto X+1$
  donne, par la formule du binôme,
  \[
    \Phi_p(X+1)=\frac{(X+1)^p-1}{X}
      =X^{p-1}+\binom{p}{1}X^{p-2}+\cdots+\binom{p}{p-1}.
  \]
  Tous les $\binom{p}{k}$, $1\leq k\leq p-1$, sont divisibles par~$p$, et le
  terme constant $\binom{p}{p-1}=p$ ne l'est pas par~$p^2$~: le \emph{critère
  d'Eisenstein} (avec le premier~$p$) s'applique, donc $\Phi_p(X+1)$, et par
  conséquent $\Phi_p(X)$, est irréductible sur $\mathbb{Q}$.

  \smallskip\noindent
  \textbf{3. Application aux corps finis (Wedderburn).}\enspace
  \index{Wedderburn (théorème de)}
  Dans la preuve du \emph{théorème de Wedderburn} (tout corps fini est
  commutatif), l'évaluation $\Phi_n(q)$ apparaît dans l'équation des classes~:
  l'entier $\Phi_n(q)$ divise $q^n-1$ et chaque $\dfrac{q^n-1}{q^d-1}$ pour
  $d\mid n$, $d<n$, ce qui force $n=1$. Le caractère \emph{entier} de
  $\Phi_n(q)$ est ici essentiel.
\end{mdframed}

\vspace{6pt}
\rubriq{Contre-exemples et pièges}

\begin{mdframed}[style=cexbar]
  \textbf{Idée fausse répandue~: « les coefficients de $\Phi_n$ sont dans
  $\{-1,0,1\}$ ».}\enspace C'est \textbf{faux}. Le premier contre-exemple est
  $n=105=3\cdot 5\cdot 7$~: le polynôme $\Phi_{105}$, de degré
  $\varphi(105)=48$, contient un coefficient égal à $-2$. (Le phénomène est lié
  au fait que $105$ est le plus petit entier produit de trois premiers impairs
  distincts.) On peut même montrer que les coefficients des $\Phi_n$ ne sont pas
  bornés lorsque $n$ varie.

  \smallskip\noindent
  \textbf{$\Phi_n$ n'est pas $X^n-1$.}\enspace On a $\Phi_n\mid X^n-1$, mais
  $\Phi_n=X^n-1$ uniquement pour $n=1$. En général $\deg\Phi_n=\varphi(n)<n$
  dès que $n\geq 2$, car $X^n-1=\prod_{d\mid n}\Phi_d$ comporte plusieurs
  facteurs.

  \smallskip\noindent
  \textbf{Racines $n$-ièmes $\neq$ racines \emph{primitives} $n$-ièmes.}\enspace
  $i$ est racine $4$-ième de l'unité \emph{et} primitive (ordre~$4$), mais
  $-1$ est racine $4$-ième \emph{sans} être primitive (ordre~$2$)~:
  $-1$ est racine de $\Phi_2$, pas de $\Phi_4$. Confondre les deux conduit à
  attribuer le mauvais degré à~$\Phi_n$.
\end{mdframed}

\vspace{6pt}
\exolabel

\begin{mdframed}[style=exobar]
  \begin{enumerate}
    \item \textit{(Eisenstein.)} Détailler la preuve de l'irréductibilité de
          $\Phi_p$ pour $p$ premier en appliquant le critère d'Eisenstein à
          $\Phi_p(X+1)$. \textit{Indication~:} $\binom{p}{k}\equiv 0\pmod p$
          pour $1\leq k\leq p-1$ et terme constant égal à~$p$.

    \item \textit{(Calcul via Möbius.)} À l'aide de
          $\Phi_n=\prod_{d\mid n}(X^d-1)^{\mu(n/d)}$, calculer $\Phi_{12}$ et
          $\Phi_{15}$. \textit{Indication~:} $\mu(1)=1$, $\mu(p)=-1$,
          $\mu(pq)=1$, $\mu(p^2)=0$.

    \item \textit{(Relation $\Phi_{2n}$.)} Montrer que pour $n\geq 3$ impair,
          $\Phi_{2n}(X)=\Phi_n(-X)$. En déduire $\Phi_6$ à partir de $\Phi_3$.
          \textit{Indication~:} $\zeta$ est racine primitive $2n$-ième ssi
          $-\zeta$ est racine primitive $n$-ième, pour $n$ impair.

    \item \textit{(Premier $p\equiv 1\bmod n$.)} Soit $p$ un premier ne divisant
          pas~$n$ tel que $p\mid \Phi_n(a)$ pour un entier~$a$. Montrer que
          l'ordre de~$\bar a$ dans $(\mathbb{Z}/p\mathbb{Z})^\times$ est
          exactement~$n$, et en déduire $p\equiv 1\pmod n$. En conclure qu'il
          existe \emph{une infinité} de premiers $\equiv 1\pmod n$ (cas
          particulier de Dirichlet).

    \item \textit{(Produit télescopique.)} Vérifier numériquement
          $X^{12}-1=\Phi_1\Phi_2\Phi_3\Phi_4\Phi_6\Phi_{12}$ et retrouver
          $\sum_{d\mid 12}\varphi(d)=12$.

    \item \textit{(Coefficients.)} Calculer $\Phi_{15}$ et constater que ses
          coefficients restent dans $\{-1,0,1\}$, puis (avec un outil de calcul
          formel) exhiber le coefficient $-2$ dans $\Phi_{105}$.
  \end{enumerate}
\end{mdframed}

\leconfooter{C.\ F.\ Gauss, \textit{Disquisitiones Arithmeticae} (1801)}
